\documentclass[11pt,a4paper]{ctexart}
% TrustBridge 中文初稿 v0.1（2026-10-24 内容版 → LaTeX）
% 编译：xelatex main_zh.tex（或 tectonic main_zh.tex）；Overleaf：编译器选 XeLaTeX
\usepackage[margin=1in]{geometry}
\usepackage{graphicx}
\usepackage{booktabs}
\usepackage{amsmath,amssymb}
\usepackage{float}
\usepackage[colorlinks=true,linkcolor=blue,citecolor=blue,urlcolor=blue]{hyperref}
\graphicspath{{../paper_figures/}}

\newcommand{\rcrf}{递归残差场}

\title{\textbf{TrustBridge：面向扩散桥医学图像翻译的\\信任图量化方法}}
\author{（作者名单待定）\\（单位）\\\small 代码与权重：\url{https://github.com/ASHDUSK/TrustBridge-}}
\date{2026 年 10 月}

\begin{document}
\maketitle

\begin{abstract}
\noindent\textbf{目的}\quad 多模态医学影像翻译模型可合成缺失模态，但生成结果的幻觉区域无法被使用者察觉，成为临床落地的主要障碍；扩散桥类模型虽已取得领先的翻译质量，其代表性工作亦将不确定性评估列为开放问题。\textbf{方法}\quad 本文提出 TrustBridge——一类无需额外训练、不依赖标注的推理期可信度量化层。该方法从自洽递归扩散桥的采样过程中提取两类被原有实现丢弃的内部信号：递归收敛残差场（刻画生成器反复修正仍无法自洽的区域）与随机桥集成分歧（刻画多次独立采样间的不一致），经分位归一化与固定权重融合为逐像素信任图，并聚合为图像级信任分与三级质控门控。\textbf{结果}\quad 在 IXI 脑 MRI、SynthRAD 盆腔 CT 与 BraTS2018 肿瘤脑共 3 个数据集、5 个翻译方向上，信任图与真实误差显著相关（AUROC 0.71--0.93）；高低信任区误差分离达 35--73 倍；消融实验证实两信号单独有效、融合最优。关键的阴性发现：在肿瘤区域（误差为瘤外的 1.6 倍）信任图完全盲（体掩码内 AUROC 0.51）——自洽性信号度量方差、无法检测偏差型幻觉，这划定了该类方法的根本边界。\textbf{结论}\quad 扩散桥的采样中间量本身即可构成有效的可信度证据；本文同时给出 5 项对照负结果，刻画了该类方法在当前数据规模下的适用边界。代码与模型已开源。

\vspace{0.5em}
\noindent\textbf{关键词：}医学图像翻译；扩散桥；不确定性量化；可信人工智能；幻觉定位
\end{abstract}

\section{引言}

多模态医学影像通过不同加权机制提供互补的组织信息，是临床诊断与放疗计划的基础。然而多协议采集显著增加扫描时间与费用，造影剂禁忌、运动伪影与回顾性研究的序列齐套需求进一步加剧了模态缺失问题。医学图像翻译旨在从已采集的源模态合成未采集的目标模态，其生成范式经历了从对抗网络~\cite{zhu2017cyclegan,isola2017pix2pix,dar2019cgan}到去噪扩散模型~\cite{ho2020ddpm,ozbey2023syndiff}、再到扩散桥~\cite{liu2023i2sb,chung2023directbridge,arslan2025selfrdb}的演进：扩散桥以前向过程直接连接源、目标两种模态分布，配合自洽递归采样，已在该任务上取得领先的翻译质量~\cite{arslan2025selfrdb}。

然而，生成模型固有的幻觉问题在医学场景中被急剧放大：合成图像中局部编造的组织细节在视觉上往往逼真，使用者无法分辨``拍摄的''与``画出来的''。一旦幻觉区域累及病灶或关键解剖结构，可能误导诊断与剂量计算。该问题已受到生成式医学影像可信化研究的广泛关注~\cite{jin2026trustworthy,fink2026airads}。值得注意的是，扩散桥的代表性工作 SelfRDB 在其讨论部分明确将``扩散桥在医学图像翻译中的不确定性评估''列为未来工作~\cite{arslan2025selfrdb}——本文正是对该开放问题的回应。

现有医学影像不确定性量化方法大多为模型级改造：MC-Dropout 需修改网络结构并重训~\cite{gal2016mcdropout}，深度集成需训练并维护多个模型~\cite{lakshminarayanan2017ensembles}，测试时增强需对输入施加人工扰动~\cite{ayhan2018tta}；针对扩散模型的集成式方法亦需多条完整采样链，且主要面向分割与分类任务~\cite{wolleb2022diffusionensembles,osuna2024microscopy}。一个自然的问题是：\textbf{在已被验证高效的扩散桥推理过程中，是否本就存在可直接提取、无需任何附加训练或标注的可信度信号？}

本文给出肯定回答。我们观察到 SelfRDB 的自洽递归采样过程产生两类被原有实现压缩或丢弃的中间量：其一，每个反向步内相邻递归估计之间的\textbf{逐像素残差}——原实现仅取其全局标量作为停机判据；其二，桥终点噪声与生成器隐变量带来的\textbf{多链采样分歧}——原文仅以``名义差异''一笔带过。我们将二者系统化为\rcrf{}与随机桥集成分歧，融合为逐像素信任图、图像级信任分与三级质控门控，并配套输入模态门控以拦截明显的用错场景。本文贡献如下：

\begin{enumerate}
  \item \textbf{方法贡献}：据我们所知，本文首次将医学图像翻译扩散桥的自洽递归中间信号系统化为无需额外训练、不依赖标注的推理期可信度证据，代价仅为集成的多次采样（可由步长跳跃部分抵消）；
  \item \textbf{实证贡献}：在 3 个数据集、5 个翻译方向上验证信任图与真实误差的强相关性（AUROC 0.71--0.93，高低信任区误差分离 35--73 倍），并给出肿瘤区信任盲区的阴性实证（见 4.5 节）；
  \item \textbf{边界贡献}：给出学习型聚合、图像级评分、自适应重采样、强度后处理、往返一致性共 5 项对照与负结果，划定方法在当前数据规模下的适用边界——逐像素信任图是该规模下唯一可靠的形态。
\end{enumerate}

\begin{figure}[htbp]
  \centering
  \includegraphics[width=\textwidth]{fig1_overview.png}
  \caption{方法总览。左：冻结的 SelfRDB 翻译模型（前向桥与含自洽递归的反向链）；右：可信度量化层——从递归过程提取残差场（信号一）与多链分歧（信号二），融合为信任图并输出图像级信任分与三级质控门控；输入端另设模态门控（虚线）。}
  \label{fig:overview}
\end{figure}

\section{相关工作}

\subsection{医学图像翻译}

早期方法以对抗训练为主：CycleGAN 借助循环一致性实现非配对翻译~\cite{zhu2017cyclegan}，pix2pix 确立配对翻译的基准范式~\cite{isola2017pix2pix}，此类方法训练不稳定且易引入伪影。扩散模型将生成建模为迭代去噪~\cite{ho2020ddpm,nichol2021iddpm}，SynDiff 以对抗式扩散实现高质量跨模态合成~\cite{ozbey2023syndiff}。扩散桥进一步将前向过程两端钉在数据分布上，直接学习源--目标迁移~\cite{liu2023i2sb,chung2023directbridge,su2023ddib}；SelfRDB 引入软先验噪声调度与自洽递归采样，取得领先性能~\cite{arslan2025selfrdb}；近期工作亦探索了桥系数的理论刻画与三维扩展~\cite{kim2024nsb,zhang2025bridge3d}，另有研究借助视觉--语言基础模型实现轻量无配对翻译~\cite{korkmaz2025i2igalip}。上述工作均以翻译质量为导向，不提供输出可信度信息。\textbf{与之相对，本文在不改动翻译模型的前提下，从其推理过程内部提取可信度证据。}

\subsection{医学影像不确定性量化}

MC-Dropout 以推断时启用 Dropout 近似贝叶斯推断~\cite{gal2016mcdropout}；深度集成通过多个独立训练模型的分歧估计不确定性~\cite{lakshminarayanan2017ensembles}；测试时增强利用输入扰动下的预测方差~\cite{ayhan2018tta}；Kendall 等对偶然与认知不确定性给出系统分类~\cite{kendall2017uncertainties}。针对扩散模型，已有工作以多采样链方差为分割与分类任务提供不确定性图~\cite{wolleb2022diffusionensembles,osuna2024microscopy}。\textbf{与上述需要模型级改造或多条完整前向链的方法不同，本文第一类信号（递归残差场）来自采样过程内部、单链即可提取；第二类信号虽同源于多链集成，但依托扩散桥仅 10 步的轻量推理，代价远低于常规扩散模型；且二者面向图像翻译而非分割/分类任务。}

\subsection{幻觉检测与可信生成}

图像级幻觉度量~\cite{tivnan2024hallucinationindex}、带不确定性输出的元数据条件合成~\cite{oliveras2026mumins}、可信生成理论框架~\cite{jin2026trustworthy}、医学 AI 输出评估框架~\cite{fink2026airads}与输出护栏机制~\cite{cardoso2023raise}共同构成生成结果可信化研究版图。\textbf{与在图像级或元数据级评估输出不同，本文提供逐像素的误差定位能力，并首次针对扩散桥翻译任务给出系统的边界分析。}

\section{方法}

\subsection{背景：自洽递归扩散桥}

给定配对训练集 $\{(x_0, y)\}$（$x_0$ 为目标模态图像，$y$ 为源模态图像），SelfRDB~\cite{arslan2025selfrdb} 在二者之间构建扩散桥。其前向过程将中间样本定义为两端的凸组合加噪：
\begin{equation}
  x_t = \mu_{x_0,t}\, x_0 + \mu_{y,t}\, y + \sigma_t\, \varepsilon,\qquad \varepsilon\sim\mathcal N(0,I),
\end{equation}
其中凸组合权重 $\mu$ 由累积扩散系数决定，噪声方差遵循软先验调度 $\sigma_t^2=\gamma\, s_t^2/(\bar s_t^2+s_t^2)$——方差沿源端单调递增，使终点成为加噪源图像 $y_\epsilon$ 而非纯噪声，反向过程因此直接完成源到目标的翻译。反向过程的每个时间步 $t$ 内，生成器 $G_\theta$ 以当前噪声图像、时间步、源图像与上一次目标估计为输入，\textbf{递归精化}目标图像估计：
\begin{equation}
  \tilde x_0^{(r+1)} = G_\theta(x_t,\ t,\ y,\ \tilde x_0^{(r)}),\qquad \tilde x_0^{(1)} = 0,
\end{equation}
直至相邻估计的相对变化低于阈值（自洽收敛），随后以收敛估计 $\tilde x_0^{*}$ 为代理执行高斯后验采样得到 $x_{t-1}$。发布配置采用 $T{=}10$ 步、每次递归上限 $R{=}2$——轻量的步数使多链集成的代价可接受，这是本文方法可行的前提。

\subsection{递归收敛残差场}

原实现将递归残差压缩为全局标量用于停机。我们保留其\textbf{逐像素}形态：在反向最后 $K$ 步（$K{=}3$，越接近数据端信号越具判别力）内收集相邻递归估计的绝对差，并在 $N$ 条采样链上平均：
\begin{equation}
  R(p) = \frac{1}{K\cdot N}\sum_{i=1}^{N}\sum_{t\in\mathcal T_K}\bigl|\tilde x_{0,i}^{(r_t+1)}(p) - \tilde x_{0,i}^{(r_t)}(p)\bigr|,
  \label{eq:rcrf}
\end{equation}
其中 $p$ 为像素坐标，$r_t$ 为步 $t$ 实际执行的递归次数。直觉上，生成器在某一空间位置反复修改仍无法自洽，表明其对该位置的结构先验薄弱，该处内容最可能为编造。该信号的提取不引入任何额外网络与前向计算，属于认知型（epistemic）不确定性证据~\cite{kendall2017uncertainties}。

\subsection{随机桥集成分歧}

扩散桥的终点噪声与生成器内部隐变量使同一输入的多次独立采样产生不同输出。我们对同一源图像独立运行 $N$ 条反向链（$N{=}4$，不同随机种子），以各链最终估计的逐像素标准差作为第二个信号：
\begin{equation}
  U(p) = \operatorname{std}_{i=1..N}\bigl(\hat x_{0,i}(p)\bigr).
\end{equation}
某位置在多链间取值不稳定，表明模型对该位置缺乏一致的生成信念。与式~\eqref{eq:rcrf} 的信号互补：残差场度量``逐步修正的稳定性''，集成分歧度量``采样的可重复性''。

\subsection{信任图融合与质控门控}

两信号分别按 99 分位归一化（抑制离群值主导）后，以固定权重融合为信任图：
\begin{equation}
  T(p) = 1 - w_1 \hat R(p) - w_2 \hat U(p),
\end{equation}
权重敏感性分析显示 $w_1$ 在 $[0.3,0.7]$ 区间内融合增益均不低于任一信号单独使用（$0.3$--$0.5$ 为性能平台，附录表~\ref{tab:weights}），本文取 $w_1{=}0.6$、$w_2{=}0.4$ 为默认。图像级信任分定义为信任图均值$\times$100，并给出三级质控门控：$\ge 75$ 为 PASS，55--75 为 REVIEW（建议人工复核），$<55$ 为 REJECT。配套输出含低信任区域叠加图与结构化报告，满足可归档、可追溯的质控要求。

\textbf{加速}。相邻后验步可合并跳跃（步长跳跃，skip=2）。该设置在信任指标持平的前提下将推理耗时降低 39\%（每片前向次数 40$\to$24，T2$\to$T1 方向 3.92 s$\to$2.38 s），已作为应用默认档位。

\subsection{输入模态门控}

信任图仅能刻画模型运行时的不确定性，无法检测``源模态用错''这类输入级错误（见~\ref{sec:limitation} 节）。为此补充一个轻量逻辑回归门控：以源图像的脑区统计特征判定其是否为任务声明的模态。在脑域内该门控对匹配与错配输入的判定正确率均不低于 0.99，可拦截明显的误用；跨解剖域（脑$\to$盆腔）泛化能力有限，作为局限如实报告。

\begin{figure}[htbp]
  \centering
  \includegraphics[width=\textwidth]{fig2_reverse_chain.png}
  \caption{真实模型的反向采样链（IXI 测试切片，官方权重）。自左至右：输入 $y$（T2）；$t{=}8,5,3$ 步内生成器对目标图像的递归估计（可见结构自粗到精逐步确立）；最终翻译结果 $\hat x_0$；真实值 $x_0$。}
  \label{fig:chain}
\end{figure}

\section{实验}

\subsection{实验设置}

\textbf{数据集}。(1)~IXI~\cite{ixi_dataset}：40 名健康受试者配对 T1/T2 脑 MRI（训 26/验 6/测 8，无受试者重叠），测试集 60 片；源/目标经世界坐标仿射对齐与跨模态细配准，强度遵循全卷均值归一与全局窗协议。(2)~SynthRAD 盆腔~\cite{thummerer2023synthrad,nyholm2018goldatlas}：A 中心 T1 加权与 C 中心 T2 加权两个子集（各 80 测试片，$128^2$），权重按中心匹配。(3)~BraTS2018~\cite{menze2015brats}：配对 FLAIR/T2 肿瘤脑 MRI，主实验 60 片、病灶分析 399 片，测试病例未经对应权重训练。

\textbf{模型}。IXI 与盆腔 CT 使用本文按官方配方自训练的权重（10 步桥、$R{=}2$、$\beta$ 自 0.1 线性至 3.0、$\gamma{=}1$、L1+对抗损失；$128^2$、30 epochs、单卡 RTX 4060）；BRATS 与盆腔 CT 另设官方发布权重对照。自训练的原因与协议见~\ref{sec:repro} 节。

\textbf{指标}。翻译质量：PSNR/SSIM（mean-norm 官方口径）与 NCC。可信度：(i)~像素级 Spearman 相关（信任图 vs 真实绝对误差）；(ii)~AUROC——以低信任为分数检测高误差像素（误差 $>90$ 分位）；(iii)~稀疏化增益 $\Delta$PSNR@$k\%$——按信任图剔除最不可信 $k\%$ 像素后剩余区域的 PSNR 提升。所有可信度指标逐片计算后取均值$\pm$标准差。

\textbf{伦理声明}：三个数据集均为公开科研数据；全部处理在本地完成，数据不出机；合成结果仅供科研演示，不用于临床决策。

\subsection{IXI 主结果}

\begin{table}[htbp]
\centering\footnotesize\setlength{\tabcolsep}{4.5pt}
\caption{IXI 主结果（测试集 60 片，自训练权重）。AUROC 为融合信任图的值；括号内为标准差。}
\label{tab:main}
\begin{tabular}{lccccccc}
\toprule
任务 & PSNR (dB) & SSIM (\%) & NCC & Spearman & AUROC（融合） & AUROC（\rcrf） & 梯度基线 \\
\midrule
T2$\to$T1 & 20.04$\pm$1.18 & 71.33 & 0.887 & $-0.806$ & \textbf{0.819$\pm$0.064} & 0.800 & 0.448 \\
T1$\to$T2 & 20.68$\pm$1.13 & 80.59 & 0.908 & $-0.902$ & \textbf{0.902$\pm$0.026} & 0.887 & 0.777 \\
\bottomrule
\end{tabular}
\end{table}

表~\ref{tab:main} 给出 IXI 双方向结果。翻译质量方面，自训练模型达到 PSNR $20.04{\pm}1.18$ / $20.68{\pm}1.13$~dB、NCC 0.887/0.908——此前误差分解表明其内容保真已达原文报告水平（线性强度校正后 31.7/32.5 dB vs 原文 31.56 dB，见~\ref{sec:repro} 节），原始口径差距源于全局强度映射残余而非结构错误。可信度方面，融合信任图与真实误差的像素级 Spearman 为 $-0.806$/$-0.902$，AUROC 达 \textbf{0.819$\pm$0.064 / 0.902$\pm$0.026}，远高于随机水平（0.5）；值得注意的是，无监督的梯度/边缘基线在 T2$\to$T1 方向完全失效（AUROC 0.448），说明信任图并非在复述``边缘处误差大''的先验，而是携带模型内部的判别信息。

\subsection{消融实验}

\begin{table}[htbp]
\centering\small
\caption{七行消融：按信任图剔除最不可信 20\% 像素后的 PSNR 增益 $\Delta$PSNR（dB）。}
\label{tab:ablation}
\begin{tabular}{lcc l}
\toprule
配置 & T2$\to$T1 & T1$\to$T2 & 说明 \\
\midrule
random & $-0.066$ & $-0.011$ & 证伪基线（随机信任图） \\
边缘/梯度 & $+0.018$ & $+1.455$ & 无监督先验基线 \\
输入差分 & $+2.022$ & $+0.385$ & 泄露真值的朴素上界参照 \\
仅 \rcrf & $+1.911$ & $+4.660$ & 创新点一单独 \\
仅集成分歧 & $+1.932$ & $+4.934$ & 创新点二单独 \\
\textbf{融合（0.6/0.4）} & $\boldsymbol{+2.136}$ & $\boldsymbol{+5.187}$ & 不低于任一单用（核心论点） \\
oracle（按真值剔除） & $+10.132$ & $+15.031$ & 上界参照 \\
\bottomrule
\end{tabular}
\end{table}

表~\ref{tab:ablation} 以稀疏化增益为口径给出消融。随机信任图增益趋零，证伪评估协议；仅递归残差场 $+1.91$/$+4.66$，仅集成分歧 $+1.93$/$+4.93$，\textbf{两信号单独均有效且量级相当；融合（$+2.14$/$+5.19$）不低于任一单用}，支持两信号互补的设计假设。以真实误差剔除的 oracle 上界为 $+10.13$/$+15.03$，表明逐像素信任图捕获了理想可剔除误差的约 21\%/35\%，为方法留有明确的改进空间。附录表~\ref{tab:weights} 进一步显示融合权重在两信号权重 0.3--0.7 区间内增益均不低于任一信号单独使用（0.3--0.5 为平台），且步长跳跃在增益持平（$+2.86$/$+5.49$）的同时节省 39\% 推理时间。

\begin{figure}[htbp]
  \centering
  \includegraphics[width=\textwidth]{fig5_ablation.png}
  \caption{消融实验柱状图（双方向）。灰色斜纹为泄露真值的参照配置，绿色斜纹为使用真值的 oracle 上界；融合配置（深蓝）在两个方向上均为最优可用方案。}
  \label{fig:ablation}
\end{figure}

\subsection{跨数据集与错配分析}
\label{sec:cross}

\begin{table}[htbp]
\centering\small
\caption{跨数据集主结果（CT 为 128\textsuperscript{2} 新切片集、80 测试片）。}
\label{tab:cross}
\begin{tabular}{llcccccc}
\toprule
数据集 & 任务 & 权重 & 片数 & PSNR (dB) & NCC & Spearman & AUROC（融合） \\
\midrule
BraTS2018 & FLAIR$\to$T2 & 官方 & 60 & \textbf{24.06$\pm$1.62} & 0.944 & $-0.897$ & \textbf{0.932$\pm$0.029} \\
盆腔 CT·A 中心（T1w） & T1$\to$CT & 官方 & 80 & 21.24$\pm$2.42 & 0.936 & $-0.673$ & 0.915$\pm$0.031 \\
盆腔 CT·C 中心（T2w） & T2$\to$CT & 官方 & 80 & 19.04$\pm$0.73 & 0.893 & $-0.626$ & 0.705$\pm$0.033 \\
盆腔 CT·A 中心（T1w） & T1$\to$CT & 自训 & 80 & \textbf{24.36$\pm$2.92} & 0.962 & $-0.808$ & 0.869$\pm$0.015 \\
盆腔 CT·C 中心（T2w） & T2$\to$CT & 自训 & 80 & \textbf{23.13$\pm$0.83} & 0.961 & $-0.766$ & 0.726$\pm$0.019 \\
\bottomrule
\end{tabular}
\end{table}

表~\ref{tab:cross} 将评估扩展至盆腔 CT 与肿瘤脑。官方权重设置下：BraTS2018 FLAIR$\to$T2 达 AUROC \textbf{0.932$\pm$0.029}（PSNR $24.06{\pm}1.62$~dB），盆腔 CT A 中心 0.915、C 中心 0.705，验证了信任图跨解剖域、跨模态对（MRI$\to$CT）的泛化性。更具启发的是自训练与官方权重的对照：自训练 CT 权重将 PSNR 提升 3--4 dB，但其信任增益趋零（$+0.36$/$-0.20$，对比官方权重的 $+4.80$/$+1.72$）。这一现象指向本文的核心论点之一：\textbf{信任信号的价值与模型--数据错配程度正相关——当模型运行于训练分布内、误差本身很小时，无需质控；而当模型被用于分布外数据（正是临床质控最需要的场景）时，信任图提供了显著的筛查能力。}换言之，官方权重在本文自备数据上的``低质量、高信任增益''并非缺陷演示，而是方法目标场景（部署后的域外运行）的真实写照。

\subsection{病灶区域的信任盲区（阴性结果）}
\label{sec:lesion}

一个核心的临床问题是：模型对``训练中未标注的病灶''是否会表现出更高的不确定性？若会，信任图即可作为无标注的病灶筛查器。我们在 BraTS2018 的 60 个含瘤切片上检验该假设（肿瘤分割标注为参照，测试病例未经对应权重训练）。

结果分两面。一方面，\textbf{误差确实集中在瘤区}：瘤内平均绝对误差 0.160，为瘤外（0.100）的 1.60 倍，且 90\% 的切片瘤内误差更高。另一方面，\textbf{信任图对瘤区完全盲}：体掩码内``低信任指示肿瘤''的 AUROC 为 0.51$\pm$0.10——精确等于随机水平；瘤内平均信任（0.772）甚至高于瘤外（0.727）。图~\ref{fig:tumor} 的案例直观展示：肿瘤（青色轮廓）被模型平滑填充为高信任的均质区域。

机理：肿瘤在 FLAIR 上表现为大块均质亮区，生成器对其\textbf{自洽地插值出同样的错误内容}——多条采样链高度一致、递归迅速收敛，故信任图给出高分；但其绝对强度与纹理是错误的（偏差型误差）。自洽性类信号度量的是\textbf{方差}（采样不一致性），而肿瘤幻觉由\textbf{偏差}主导——方差型信号对偏差型误差在原理上不可见。该阴性结果与~
\ref{sec:cross} 节的错配分析共同构成完整图景：信任图能定位``模型犹豫''的随机性误差，无法定位``模型自信地错''的系统性偏差；后者恰是临床最危险的幻觉类型，需要外部先验（而非采样统计）来检测。

\begin{figure}[htbp]
\centering
\includegraphics[width=\textwidth]{fig6_tumor.png}
\caption{病灶分析（修正版，BraTS2018，60 个含瘤切片，体掩码内口径）。(a) 瘤内/体（非瘤）逐片平均信任值；(b) 逐片平均绝对误差——瘤内为瘤外的 1.60 倍；(c) 体掩码内低信任指示肿瘤的 AUROC 分布——均值 0.51，等于随机水平；(d--f) 案例切片：FLAIR（瘤区亮）、信任图+瘤轮廓（瘤区为高信任）、翻译结果。信任图对瘤区完全盲：模型在均质亮区自洽地插值出错误内容（偏差型幻觉），方差型信号不可见。}
\label{fig:tumor}
\end{figure}

\begin{table}[htbp]
\centering\small
\caption{病灶分析（修正版，BraTS2018，60 个含瘤切片，体掩码内口径）。(a) 瘤内/体（非瘤）逐片平均信任值；(b) 逐片平均绝对误差——瘤内为瘤外的 1.60 倍；(c) 体掩码内低信任指示肿瘤的 AUROC 分布——均值 0.51，等于随机水平；(d--f) 案例切片：FLAIR、信任图+瘤轮廓（瘤区为高信任）、翻译结果。}
\label{tab:lesion}
\begin{tabular}{lcc}
\toprule
指标 & 瘤内 & 体（非瘤）\\
\midrule
平均信任值 & 0.772 & 0.727 \\
平均绝对误差 & 0.160 & 0.100 \\
\midrule
低信任指示肿瘤 AUROC（体掩码内） & \multicolumn{2}{c}{0.51$\pm$0.10（随机水平）} \\
误差比（瘤内/瘤外） & \multicolumn{2}{c}{1.60$\times$（90\% 切片瘤内更高）} \\
\bottomrule
\end{tabular}
\end{table}

\subsection{校准分析}

\begin{figure}[htbp]
  \centering
  \includegraphics[width=\textwidth]{fig4_reliability.png}
  \caption{可靠性分析（每方向 80 片）：按信任值十分箱后各箱的平均绝对误差。两方向均近单调递减，最高/最低信任箱误差分离比分别为 35.3$\times$ 与 73.2$\times$。}
  \label{fig:calibration}
\end{figure}

除排序能力外，我们检验信任分是否具有校准含义。将每片像素按信任值十分箱（图~\ref{fig:calibration}），各箱平均误差在两个方向上均近单调递减（单调性 0.889），最高信任箱与最低信任箱的平均误差分离比达 \textbf{35.3$\times$（0.72$\to$0.020）与 73.2$\times$（1.07$\to$0.015）}。这说明信任图不是任意排序，其数值幅度携带``可信程度''的绝对信息，支撑图像级信任分与三级门控的语义可用性。

\begin{table}[htbp]
\centering\small
\caption{校准分析（每方向 80 片）。}
\label{tab:calibration}
\begin{tabular}{lccc}
\toprule
方向 & 单调性 & 最高/最低信任箱误差分离比 & Spearman（信任，误差）\\
\midrule
T2$\to$T1 & 0.889 & 35.3$\times$（0.72 $\to$ 0.020） & $-0.806$ \\
T1$\to$T2 & 0.889 & 73.2$\times$（1.07 $\to$ 0.015） & $-0.903$ \\
\bottomrule
\end{tabular}
\end{table}

\section{讨论}

\subsection{五项对照与负结果：方法的适用边界}

我们系统评估了六类``看起来应有效''的增强或假设，均未取得增益或被证伪。这些负结果共同刻画了方法边界：(i)~\textbf{肿瘤信任盲区}（\ref{sec:lesion} 节）：瘤区误差 1.6 倍于瘤外而信任图 AUROC 仅 0.51——方差型信号对偏差主导的误差不可见，是临床意义上最重要的边界；(ii)~\textbf{学习型融合}（梯度提升树，块级划分）在两方向上一升一降（0.830/0.869 vs 固定 0.819/0.902），不稳定；(iii)~\textbf{图像级 QC 多特征模型}测试集全面低于简单基线（0.323/0.325 vs 0.381/0.401），跨方向池化时方向效应互相抵消；(iv)~\textbf{自适应采样与重采样}（块级早停、预算再分配、择优保留）与均匀基线差异在 $-0.05$~dB 以内；离线模拟表明链均值已饱和；RePaint 式区域重采样仅获 $+0.10$~dB 的初步增益——采样策略层面无实质性免费增益；原(iii)项并入本项，，且离线模拟表明链均值已饱和（8 链 vs 4 链仅 $+0.02$~dB）——采样策略层面无免费增益；(v)~\textbf{强度标准化后处理}（分位数匹配/仿射矩匹配）使 PSNR 与 SSIM 双降（$-0.4$至$-0.9$~dB），说明扩散桥的强度偏差并非全局可校正项；(vi)~\textbf{往返一致性信号}单独 AUROC 仅 0.503，加入后反而稀释融合性能。综合结论：\textbf{在当前数据规模（数百至数千训练切片）下，逐像素、固定融合的信任图是唯一可靠的形态}。这些负结果与\ref{sec:lesion} 节的肿瘤盲区共同构成对自洽性可信度信号的完整边界刻画：\textbf{它们能定位方差型（随机性）误差，不能定位偏差型（系统性）幻觉。}

\subsection{与现有不确定性方法的关系}

本文第二信号在形式上属于深度集成~\cite{lakshminarayanan2017ensembles}的推理期变体，但其成立依赖扩散桥的两个特性：10 步轻量推理使多链代价可接受；终点为加噪源图像使各链共享强先验、分歧集中于真正不确定的区域。第一信号（递归残差场）则据我们所知为扩散桥特有——常规 DDPM 无递归结构，不存在对应中间量。二者的互补性由消融（表~\ref{tab:ablation}）与权重敏感性（附录表~\ref{tab:weights}）共同支持。肿瘤盲区的发现也与不确定性文献的经典结论呼应：不确定性量化捕获的是方差，系统性偏差需要分布外检测等外部机制——本文的实证为此提供了医学翻译任务上的直接证据。

\subsection{复现性说明}
\label{sec:repro}

官方 IXI 权重在本文自备数据上的重建质量约为原文报告的一半，经系统排查（对齐、强度窗、裁剪、方向等），可归因于原文未完全公开的数据制备细节，该现象与社区报告一致。为此本文全部定量结论基于按官方配方自训练的闭环流程（数据制备、训练、评估全部开源），官方权重仅用于跨数据集对照——这恰构成~\ref{sec:cross} 节``错配场景''的实验素材。

\subsection{局限}
\label{sec:limitation}

(i)~信任图标记的是\textbf{不可靠}区域而非\textbf{错误}区域：自信地合成的幻觉对信任图不可见；(ii)~图像级 QC 未解：简单聚合与学习型聚合均不可靠，信任分的门控语义依赖像素级信任图的校准性；(iii)~输入门控跨解剖域泛化失败（脑域正确率不低于 0.99，脑$\to$盆腔错配仅拦截 20/90）；(iv)~病灶指示结论依赖``测试病例未经训练''，在训练已见病灶的设置下信号失效；(v)~全部评估为 2D 切片级，3D 上下文未利用。

\section{结论与未来工作}

本文从自洽递归扩散桥的推理过程中提取两类被丢弃的内部信号，构建了无需额外训练、不依赖标注的医学图像翻译可信度量化层。三个数据集、五个方向的系统评估表明：信任图与真实误差强相关且具备校准含义；同时，肿瘤盲区的阴性实证划定了其根本边界——\textbf{自洽性类信号度量采样方差，能够定位``模型犹豫''的随机性误差，但无法检测``模型自信地错''的偏差型幻觉}。这一边界对未来工作具有直接指导意义：可信度证据必须与外部先验（解剖约束、检测器、放疗物理约束）结合，才能覆盖自信的系统性幻觉。近期将推进：扩大受试者规模以重启可学习聚合、跨中心 CT 序列自动匹配、RePaint 式区域重采样（已有 $+0.10$~dB 的初步阳性信号）的深入验证，以及在真实临床工作流中的前瞻性质控评估。

代码、模型权重与全部逐片实验数据已开源：\url{https://github.com/ASHDUSK/TrustBridge-}。

\section*{致谢}
感谢 SelfRDB 作者（Bilkent ICON Lab）以 MIT 协议开源代码与权重，感谢 IXI、SynthRAD 与 BraTS 各数据集联盟的公开数据。本研究为学生科研演示，输出不得用于临床决策。

\begin{thebibliography}{99}\small

\bibitem{arslan2025selfrdb} Arslan F, Kabas B, Dalmaz O, Ozbey M, {\c C}ukur T. Self-consistent recursive diffusion bridge for medical image translation. arXiv:2405.06789 (Medical Image Analysis), 2024/2025.
\bibitem{zhu2017cyclegan} Zhu J Y, Park T, Isola P, Efros A A. Unpaired image-to-image translation using cycle-consistent adversarial networks. ICCV, 2017.
\bibitem{isola2017pix2pix} Isola P, Zhu J Y, Zhou T, Efros A A. Image-to-image translation with conditional adversarial networks. CVPR, 2017.
\bibitem{dar2019cgan} Dar S U H, Yurt M, Karacan L, Erdem A, Erdem E, {\c C}ukur T. Image synthesis in multi-contrast MRI with conditional generative adversarial networks. IEEE Trans Med Imaging, 2019, 38(10): 2375--2388.
\bibitem{ho2020ddpm} Ho J, Jain A, Abbeel P. Denoising diffusion probabilistic models. NeurIPS, 2020.
\bibitem{nichol2021iddpm} Nichol A Q, Dhariwal P. Improved denoising diffusion probabilistic models. ICML, 2021.
\bibitem{liu2023i2sb} Liu G H, Vahdat A, Huang D A, Theodorou E A, Nie W, Anandkumar A. I2SB: Image-to-image Schr{\"o}dinger bridge. ICML, 2023.
\bibitem{chung2023directbridge} Chung H, Kim J, Ye J C. Direct diffusion bridge using data consistency for inverse problems. arXiv:2305.19809, 2023.
\bibitem{su2023ddib} Su X, Song J, Meng C, Ermon S. Dual diffusion implicit bridges for image-to-image translation. ICCV, 2023.
\bibitem{kim2024nsb} Kim B, Kwon G, Kim K, Ye J C. Unpaired image-to-image translation via neural Schr{\"o}dinger bridge. ICLR, 2024.
\bibitem{zhang2025bridge3d} Zhang 等. 三维扩散桥系数理论相关工作. 2025.（核实条目，见 docs/15）
\bibitem{korkmaz2025i2igalip} Korkmaz Y, Patel V M. I2I-Galip: Unsupervised medical image translation using generative adversarial CLIP. MIDL, 2025.
\bibitem{ozbey2023syndiff} {\O}zbey M, Dalmaz O, Dar S U H, et al. Unsupervised medical image translation with adversarial diffusion models. IEEE Trans Med Imaging, 2023, 42(12): 3524--3539.
\bibitem{gal2016mcdropout} Gal Y, Ghahramani Z. Dropout as a Bayesian approximation: Representing model uncertainty in deep learning. ICML, 2016.
\bibitem{lakshminarayanan2017ensembles} Lakshminarayanan B, Pritzel A, Blundell C. Simple and scalable predictive uncertainty estimation using deep ensembles. NeurIPS, 2017.
\bibitem{ayhan2018tta} Ayhan M S, Berens P. Test-time data augmentation for estimation of heteroscedastic aleatoric uncertainty. 2018.
\bibitem{kendall2017uncertainties} Kendall A, Gal Y. What uncertainties do we need in Bayesian deep learning for computer vision? NeurIPS, 2017.
\bibitem{wolleb2022diffusionensembles} Wolleb J, Sandk{\"u}hler R, Bieder F, et al. Diffusion models for medical anomaly detection. MICCAI, 2022.
\bibitem{osuna2024microscopy} Osuna kurachi 等. 扩散集成不确定性（显微成像）. 2024.（核实条目，见 docs/15）
\bibitem{tivnan2024hallucinationindex} Tivnan D, et al. Benchmarking hallucination in AI-generated medical images: The Hallucination Index. 2024.
\bibitem{oliveras2026mumins} Oliveras R, et al. MUMINS: Metadata-conditioned synthesis with uncertainty maps. 2026.
\bibitem{jin2026trustworthy} Jin 等. Trustworthy generative AI framework. 2026.
\bibitem{fink2026airads} Fink 等. 医学 AI 输出评估框架（AI-RADS 类）. 2026.（作者字段待补全）
\bibitem{cardoso2023raise} Cardoso A, et al. RAISE: 输出护栏机制. 2023.
\bibitem{ixi_dataset} IXI Dataset. \url{https://brain-development.org/ixi-dataset/}.
\bibitem{menze2015brats} Menze B H, et al. The multimodal brain tumor image segmentation benchmark (BraTS). IEEE Trans Med Imaging, 2015, 34(10): 1993--2024.
\bibitem{thummerer2023synthrad} Thummerer A, et al. SynthRAD2023 grand challenge. 2023.
\bibitem{nyholm2018goldatlas} Nyholm T, et al. MR and CT data with multiobserver delineations of organs in the pelvic area---part of the Gold Atlas project. Med Phys, 2018, 45(3): 1295--1300.

\end{thebibliography}

\appendix

\section{融合权重敏感性与步长跳跃}

\begin{table}[htbp]
\centering\small
\caption{融合权重敏感性（$w_1$ 为递归残差场权重）与步长跳跃加速（IXI 测试集 60 片，$\Delta$PSNR@20\%，dB）。}
\label{tab:weights}
\begin{tabular}{lcc}
\toprule
$w_1$（$w_2=1-w_1$） & T2$\to$T1 & T1$\to$T2 \\
\midrule
1.0（纯 \rcrf） & 2.403 & 4.728 \\
0.7 & 2.712 & 5.248 \\
0.5 & 2.811 & 5.439 \\
\textbf{0.3} & \textbf{2.841} & \textbf{5.472} \\
0.0（纯集成分歧） & 2.448 & 4.900 \\
\midrule
\multicolumn{3}{l}{\emph{步长跳跃 skip=2（融合权重 0.6/0.4）：}} \\
\quad skip=1（逐步，40 次前向） & 2.771（3.92 s/片） & 5.354（3.10 s/片） \\
\quad skip=2（24 次前向，$-39\%$ 耗时） & 2.859（2.38 s/片） & 5.494（1.89 s/片） \\
\bottomrule
\end{tabular}
\end{table}

两信号任一端点单独使用均损失 0.4--0.6 dB，固定融合在 $w_1\in[0.3,0.5]$ 平台内稳健；步长跳跃在节省 39\% 时间的同时指标处于噪声内微升，已作为应用默认档位。

\end{document}
